% ============================================================
%  CAPÍTULO 1: INTRODUCCIÓN
% ============================================================
\chapter{Introducción}
\label{cap:introduccion}

% ------------------------------------------------------------
\section{Motivación y contexto del problema}
% ------------------------------------------------------------

El Sistema Nacional de Salud (SNS) español atraviesa una crisis estructural cuya
manifestación más visible son las listas de espera. Según datos del
Ministerio de Sanidad, a finales de 2023 más de 800.000 pacientes esperaban una
intervención quirúrgica, con tiempos medios que en algunas comunidades autónomas
superaban los 150 días para consultas con especialistas. Esta situación no es
coyuntural: desde 2016 se observa una tendencia sostenida al alza, interrumpida
únicamente por la pandemia de la COVID-19, período en el que se paralizó la 
actividad programada y que después se disparó.

La persistencia y magnitud del problema de las listas de espera plantea
dudas fundamentales sobre la sostenibilidad del modelo sanitario español
y, de forma más específica, sobre las respuestas de los hogares ante el 
deterioro percibido del servicio público. En un sistema de 
cobertura universal y gratuita (en cuanto al punto de uso), ¿qué factores determinan que
cada vez una proporción mayor de la población opte por contratar un seguro sanitario
privado complementario? ¿Constituyen los tiempos de espera un factor de expulsión
 calculable empíricamente? ¿Y, de ser así, cómo evoluciona este mecanismo en el tiempo?

Este trabajo aborda estas cuestiones desde la economia de la salud y los 
modelos de datos de panel, combinando un marco teórico riguroso con
evidencia empírica original para España.

% ------------------------------------------------------------
\section{La teoría del precio sombra aplicada a la sanidad pública}
% ------------------------------------------------------------
Este trabajo parte de la premisa de que la sanidad publica, a pesar de ser gratuita en el punto de uso,
tiene unos costes implicitos para el usuario de esta, siendo el tiempo de espera el principal de ellos.
Semanas o meses de incertidumbre sobre diagnosticos, dolores no tratados, incapacidad laboral o deteriro de la 
calidad de vida son costes temporales no monetarios quye pueden estudiarse a traves del concepto de precio sombra.

El precio sombra del tiempo de espera se puede entender como el valor no monetario que el hogar asigna
al periodo que transcurre desde que se pide la cita con el especialista hasta que finalmente acude a 
la consulta. Para un trabajador con un salario alto cada dia de espera tiene un coste de oportunidad mayor (salario 
renunciado o productividad perdida) o para un paciente con un problema de salud progresivo, cada dia sin 
tratamiento puede suponer un deterioro de su estado de salud y calidad de vida. En ambos casos, cuando el precio
sombra de la sanidad publica aumenta (por ejemplo, por un incremento de los tiempos de espera), el precio 
relativo de los seguros privados disminuye, incentivando su demanda. 

Esta perspectiva conecta directamente con la propuesta de \textcite{Grossman1972}
sobre la demanda de salud como capital humano y con la literatura posterior sobre
elección entre sectores sanitarios \parencite{Propper2000, Besley1999}. La
aportación específica de este trabajo consiste en modelizar el precio sombra
mediante los dias de espera para la consulta con especialistas y contrastar
empíricamente su efecto sobre el gasto en seguros privados de los hogares
españoles.

% ------------------------------------------------------------
\section{Pregunta de investigación e hipótesis}
% ------------------------------------------------------------

La pregunta central que guía este trabajo es:

\begin{quote}
\textit{¿Cuáles son los determinantes del gasto de los hogares españoles en
seguros sanitarios privados, y qué papel desempeñan los tiempos de espera del
sistema público en esta decisión?}
\end{quote}

A partir del marco teórico, se formulan las siguientes hipótesis:

\begin{description}
  \item[\textbf{H1} (Efecto renta):] El gasto en seguros sanitarios crece con la renta disponible 
    como un bien normal o superior.
  
  \item[\textbf{H2} (Efecto precio sombra):] Si aumentan los tiempos de espera del sistema 
    publico, el gasto en seguros privados se incrementa ya que se reduce el precio relativo de estos.

  \item[\textbf{H3} (Efecto diferido):] la respuesta de los hogares a los tiempos de 
    espera no es inmediata, sino que hay un retraso temporal de uno o dos años 
    debido a la necesidad de que los hogares perciban el deterioro y busquen alternativas.

  \item[\textbf{H4} (Efecto umbral):] Existe un umbral crítico de tiempo de espera
    (estimado en 100 días) a partir del cual la respuesta de los hogares
    se intensifica y por tanto el efecto no es lineal.

  \item[\textbf{H5} (Persistencia del gasto):] El gasto en seguros privados es persistente en el
    tiempo debido a la naturaleza de los contratos, de modo que el consumo pasado influye
    positivamente en el presente. 
\end{description}

% ------------------------------------------------------------
\section{Originalidad y contribución del trabajo}
% ------------------------------------------------------------

Este Trabajo de Fin de Grado realiza varias contribuciones originales:

\textbf{Primero}, construye un panel de datos autonomico (17 CCAA,
2016--2024) que combina información del INE (gasto en seguros, PIB per cápita,
demografía) y del Ministerio de Sanidad (tiempos de espera) puediendo asi estudiar la evolución
temporal y regional de las variables clave.

\textbf{Segundo}, introduce la Ley de Little de la teoría de colas como
fundamento para la elección de los días de espera, y no el número de pacientes
en lista, como proxy de saturación del sistema público percibida por el
hogar.

\textbf{Tercero}, estima modelos econométricos que capturan el efecto difereido
 de los tiempos de espera sobre el gasto privado utilizando retardos de uno y dos años
 cuya significatividad se contrasta mediante tests de Wald.

\textbf{Cuarto}, explora el hecho de que el efecto de los tiempos de espera pueda no ser
lineal y que se dispare a partir de cierto umbral de dias.

\textbf{Quinto}, analiza la persistencia dinamica del gasto en seguros privados, coherente con la 
naturaleza de los contratos y la fidelizacion de los consumidores.

% ------------------------------------------------------------
\section{Estructura del trabajo}
% ------------------------------------------------------------

El trabajo se organiza en nueve capítulos:

\begin{itemize}
  \item El \textbf{Capítulo~\ref{cap:contexto}} presenta el marco institucional del
    SNS español, la evolución histórica de las listas de espera y el papel del
    sector privado.
  \item El \textbf{Capítulo~\ref{cap:teoria}} desarrolla el marco teórico: modelo
    de decisión del hogar, Ley de Little y precio sombra del tiempo de espera.
  \item El \textbf{Capítulo~\ref{cap:datos}} describe las fuentes de datos, las
    variables y los estadísticos descriptivos del panel.
  \item El \textbf{Capítulo~\ref{cap:metodologia}} desarrolla la metodología
    econométrica: efectos fijos/aleatorios, tests de cointegración, errores
    estándar agrupados y modelo dinámico.
  \item El \textbf{Capítulo~\ref{cap:resultados}} presenta los resultados de las
    estimaciones y su interpretación económica.
  \item El \textbf{Capítulo~\ref{cap:discusion}} discute las implicaciones de los
    hallazgos para la política sanitaria.
  \item El \textbf{Capítulo~\ref{cap:limitaciones}} reconoce las limitaciones del
    trabajo.
  \item El \textbf{Capítulo~\ref{cap:conclusiones}} sintetiza las conclusiones y
    propone líneas de investigación futura.
\end{itemize}
